\chapter{Experiments and Results}
\label{ch:results}

\section{Experimental setup}

All results below are produced by the pipeline described in
Chapter~\ref{ch:method}: a 3D~U-Net trained on the 249 high-grade-glioma training
subjects and evaluated on the 44 held-out validation subjects, with the four MRI
modalities as input and the four-class segmentation as output. Metrics are
computed on the validation split using the best checkpoint, which was obtained at
epoch~34; as Figure~\ref{fig:curves} shows, the validation mean Dice had
plateaued by that point, so this checkpoint is representative of the model's
converged performance for this configuration.

\begin{figure}[t]
\centering
\includegraphics[width=0.9\textwidth]{sample_data.png}
\caption{A representative training subject: the four input modalities
(T1, T1Gd, T2, FLAIR) and the expert segmentation overlaid on FLAIR. The mask
distinguishes the necrotic/non-enhancing core, peritumoral edema and enhancing
tumour. The four modalities carry complementary information, which motivates using
all of them as input.}
\label{fig:sample}
\end{figure}

\section{Training behaviour}

Figure~\ref{fig:curves} shows the training and validation loss together with the
validation mean Dice over the course of training. The loss decreases smoothly for
both splits and the validation Dice rises and then plateaus, at which point early
stopping selects the best checkpoint. The small gap between training and
validation loss indicates that the light augmentation and dropout keep
overfitting under control on this modestly sized dataset.

\begin{figure}[t]
\centering
\includegraphics[width=0.95\textwidth]{training_curves.png}
\caption{Training dynamics. Left: training and validation loss per epoch.
Right: validation mean Dice per epoch. The validation Dice rises steeply over the
first ten epochs and plateaus around 0.77 by epoch~20; the small, stable gap
between training and validation loss indicates limited overfitting.}
\label{fig:curves}
\end{figure}

\section{Quantitative results}

Table~\ref{tab:results} reports the Dice similarity coefficient and the
95th-percentile Hausdorff distance for each tumour sub-region on the validation
split, together with the foreground mean. Dice is reported in $[0,1]$ (higher is
better) and HD95 in millimetres (lower is better). The model reaches a foreground
mean Dice of 0.71 with a mean HD95 of 9.1~mm. Peritumoral edema and enhancing
tumour are segmented most accurately (Dice 0.77 and 0.73), while the
necrotic/non-enhancing core is the hardest sub-region (Dice 0.61), consistent with
its small size and ambiguous, enhancement-defined boundary. These figures sit in
the expected range for a compact single-GPU 3D~U-Net on BraTS.

\begin{table}[t]
\centering
\caption{Segmentation performance on the 44-subject HGG validation split,
using the best checkpoint (epoch~34). Dice higher is better; HD95 lower is better.}
\label{tab:results}
\begin{tabular}{@{}lcc@{}}
\toprule
Sub-region & Dice $\uparrow$ & HD95 (mm) $\downarrow$ \\
\midrule
Necrotic / non-enhancing core (NCR/NET) & 0.614 & 8.7 \\
Peritumoral edema (ED)                  & 0.767 & 12.0 \\
Enhancing tumour (ET)                   & 0.735 & 6.4 \\
\midrule
Foreground mean                         & 0.705 & 9.1 \\
\bottomrule
\end{tabular}
\end{table}

\section{Qualitative results}

Figure~\ref{fig:overlay} compares the model's prediction with the ground truth on
a validation case, both overlaid on the FLAIR image. On this subject the
prediction matches the ground truth closely across all three sub-regions: the
enhancing tumour (blue), the necrotic/non-enhancing core (red) and the surrounding
edema (green) are each recovered with good spatial agreement. Across the full
validation set the model is most reliable on the edema and enhancing tumour and
weakest on the small, ambiguous necrotic core, consistent with
Table~\ref{tab:results}.

\begin{figure}[t]
\centering
\includegraphics[width=0.95\textwidth]{overlay.png}
\caption{Qualitative segmentation on a validation subject: input MRI (FLAIR),
ground-truth mask and model prediction (green: edema, red: necrotic/non-enhancing
core, blue: enhancing tumour). The prediction matches the ground truth closely
across all three sub-regions.}
\label{fig:overlay}
\end{figure}

\section{Per-subject variability}
\label{sec:variability}

The aggregate metrics in Table~\ref{tab:results} average over all 44 validation
subjects, but segmentation difficulty varies substantially from case to case.
Figure~\ref{fig:dice-dist} shows the distribution of per-subject foreground mean
Dice: scores range from 0.22 to 0.94, with a mean of 0.71 (matching the aggregate
in Table~\ref{tab:results}), a median of 0.76 and a standard deviation of 0.17.
The necrotic/non-enhancing core has the widest spread and includes several
subjects with near-zero Dice --- typically cases where the true core is very
small or entirely absent, so a modest false positive or a missed detection
dominates the ratio. Edema and enhancing tumour are comparatively stable across
subjects, consistent with their higher aggregate Dice in Table~\ref{tab:results}.

\begin{figure}[t]
\centering
\includegraphics[width=0.75\textwidth]{dice_distribution.png}
\caption{Per-subject foreground Dice distribution across the 44 validation
subjects, split by class. Box: interquartile range; orange line: median; green
dashed line: mean; dots: individual subjects.}
\label{fig:dice-dist}
\end{figure}

\section{Qualitative results across performance levels}

Figure~\ref{fig:overlay} showed a single representative case. To give a fuller
picture, Figure~\ref{fig:cases} shows the \emph{best}, \emph{typical} (median)
and \emph{worst} validation subjects, selected by foreground mean Dice from the
distribution in Figure~\ref{fig:dice-dist} rather than chosen by hand. In the
best case (Dice 0.94, subject BraTS20\_Training\_192) the model recovers the
classic concentric ring structure of glioblastoma --- necrotic core, enhancing
rim, surrounding edema --- almost exactly. In the typical case (Dice 0.77,
subject BraTS20\_Training\_013) the overall structure is correct with minor
boundary disagreement. In the worst case (Dice 0.22, subject
BraTS20\_Training\_141) the model detects the correct general location via the
edema but fails to resolve the internal necrotic core and enhancing tumour at
all, illustrating the sub-structure failure mode discussed further in
Section~\ref{sec:failure-cases}.

\begin{figure}[t]
\centering
\includegraphics[width=0.95\textwidth]{case_best.png}\\[4pt]
\includegraphics[width=0.95\textwidth]{case_typical.png}\\[4pt]
\includegraphics[width=0.95\textwidth]{case_worst.png}
\caption{Best, typical and worst validation cases by foreground mean Dice
(top to bottom), each showing MRI (FLAIR), ground truth and prediction at the
subject's peak-tumour slice.}
\label{fig:cases}
\end{figure}

\section{Comparison with published BraTS results}
\label{sec:literature-comparison}

The evaluation elsewhere in this chapter uses the three disjoint BraTS labels
(necrotic/non-enhancing core, edema, enhancing tumour) defined in
Section~\ref{sec:preproc}. The official BraTS challenge instead scores three
overlapping composite regions: the whole tumour (WT, all abnormal tissue), the
tumour core (TC, necrotic core plus enhancing tumour) and the enhancing tumour
(ET) alone. To allow a direct, like-for-like comparison with published results,
we additionally compute these composite regions from the same predictions.
Table~\ref{tab:literature} places our result alongside the top two entries in
the BraTS~2020 segmentation challenge.

\begin{table}[t]
\centering
\caption{Composite-region Dice compared with the top two BraTS~2020 challenge
entries. WT: whole tumour; TC: tumour core; ET: enhancing tumour.}
\label{tab:literature}
\begin{tabular}{@{}lccc@{}}
\toprule
Method & WT & TC & ET \\
\midrule
nnU-Net \citep{isensee2021nnunet} (1st place, test set)     & 0.890 & 0.851 & 0.820 \\
H2NF-Net \citep{jia2020h2nfnet} (2nd place, validation set) & 0.913 & 0.855 & 0.788 \\
\midrule
This work (compact 3D~U-Net, validation split)              & 0.884 & 0.817 & 0.735 \\
\bottomrule
\end{tabular}
\end{table}

The comparison should be read with its context in mind: the challenge entries
are ensembles trained with extensive compute budgets, challenge-specific
post-processing and, in nnU-Net's case, an automated pipeline-configuration
search, whereas this model is a single, untuned 3D~U-Net trained for 34 epochs
on one 8~GB consumer GPU. Even so, whole-tumour localisation (WT, 0.884) sits
within a few points of both challenge entries, suggesting the model has
learned to find the tumour reliably. The gap widens for the tumour core and,
most visibly, the enhancing tumour (0.735 versus 0.79--0.82), which matches the
disjoint-label finding in Table~\ref{tab:results} that the smaller, more
detailed sub-regions are where this model has the most room to improve with
additional training or capacity.
